\documentclass[12pt]{article}
\usepackage[utf8]{inputenc}
\usepackage[T1]{fontenc}
\usepackage{hyperref}
\title{Fali Muchella Language Resources for K-12 Curriculum Development}
\author{Compiled by Claude Sonnet 5.5 (claude-sonnet-5-5) for Andr\'as Kornai}
\date{}
\begin{document}
\maketitle

\section{Introduction}
Fali Muchella (ISO 639-3 fli; Glottocode fali1285) is spoken in Adamawa State, Nigeria (Mubi area). It should not be confused with the Fali languages of northern Cameroon. Population figures in the sources disagree (about 20,000 to 38,000 for related Fali groups). Search turned up no dictionary, grammar or teaching materials specific to this language.

\section{Primary Educational Resources}
None found.

\section{Linguistic and Reference Materials}
None specific to Fali Muchella were found. A German-language book on the Fali of northern Cameroon appeared in results but concerns a different language and is deliberately not cited.

\section{Religious and Community Resources}
\begin{itemize}
\item Global Recordings Network, Fali (Madzarin) page, listing Fali-language gospel audio-visual lessons and audio programmes: \url{https://globalrecordings.net/en/language/9754}
\item Joshua Project, Vimtim Fali in Nigeria (related Fali group; scripture portions and audio reported): \url{https://joshuaproject.net/people_groups/11181}
\end{itemize}

\section{Digital and Audio Resources}
Audio materials are limited to the Global Recordings Network items above.

\section{Curriculum Development Context}
The language is stable and used by all in the community, but no educational or reference materials were found. Curriculum work would begin from primary documentation.
\end{document}